% =====================================================================
%  5. DISCUSSION
% =====================================================================
\section{Discussion}\label{sec:discussion}

\subsection{Ce que l’étude établit}

Un protocole se juge à ce qu’il rend impossible. Sur les trois comparatifs
comme sur l’extension rétro-active, l’étude établit que la classe de
défauts « contrat non écrit » — valeur-seuil enfouie, structure déclarée de
mémoire, comportement dépendant de l’hôte, analyse de libellés localisés —
est systématiquement convertie en contrat nommé, documenté, testé et
journalisé. Les règles ont opéré sans exception sur un périmètre réel et
complet~: 42 modules cartographiés, 44 spécifications écrites avant le
code, dix commits tous précédés de la vérification complète, sept vagues
closes par une sauvegarde systématique. L’oracle a joué son rôle jusque
dans la révélation de défauts invisibles sur la machine de développement
(le caractère interdit dépendant de l’hôte,
section~\ref{sec:caractere}~; le verrouillage de journal révélé par la CI
Windows, section~\ref{sec:resultats}). La parité est démontrable par des
tiers — empreintes reproductibles, quatorze familles auditées, écarts
documentés — et la structure s’est améliorée sans gonflement~: à volume de
code comparable, la défense double (228~$\to$~437 tests, avertissements
tolérés~$\to$~zéro) et le plus gros fichier applicatif passe de 1\,075 à
399~lignes.

La démonstration valide aussi, indirectement, la \emph{gouvernance}. Le
journal des décisions — quinze entrées en cinq sessions, chacune avec ses
alternatives rejetées — a servi de mémoire d’arbitrage effectivement
consultée~: aucune révision silencieuse, aucun accommodement improvisé, y
compris sous la pression d’une réinitialisation d’environnement en cours
de projet. L’anti-enfermement n’est pas théorique~: l’historique, les
contrats et l’état sont récupérables par un agent neuf depuis les seuls
artefacts du dépôt. Et la séquence v3.0.0~$\to$~v3.0.2 (deux échecs de CI
honorés puis corrigés en patchs à contrat stable) montre que la
discipline tient \emph{dans l’échec}, ce qui est sa véritable épreuve.

\subsection{Implications pour le génie logiciel assisté par IA}

Pour l’AISE, l’étude suggère un déplacement du regard~: à capacité du
modèle constante, la qualité du résultat agent-first dépend du
\emph{processus} qui encadre l’agent. Les modes d’échec documentés de la
littérature — code plausible mais inexact, dérive silencieuse hors
périmètre, dépendances hallucinées \cite{hou2024} — ne sont pas que des
propriétés des modèles~; ce sont des propriétés de \emph{couples}
modèle-processus, et le second terme est actionnable dès aujourd’hui. Les
trois mécanismes que l’étude a éprouvés se transposent à tout
outillage~: (i)~la spécification écrite avant le code, qui force
l’intention à précéder l’exécution et rend chaque évolution
\emph{explicable}~; (ii)~le vérifie-puis-écrit (R5), qui transforme la
confiance dans le modèle en vérification de l’artefact — l’agent peut
halluciner, le commit reste propre~; (iii)~la cristallisation à
pré-versions explicites, qui distingue « vérifié par l’oracle » de
« validé sur le terrain » et interdit la confusion des niveaux de preuve.
La reverse-spec, enfin, adresse le problème de l’\emph{embarquement} de
patrimoine~: face à un legs sans documentation, l’agent ne repart pas de
zéro — il extrait d’abord le contrat du code relu, puis le fait valider.

\subsection{Coûts et frictions observés}

L’honnêteté impose de chiffrer la contrepartie. La discipline a un coût
d’écriture~: 44 spécifications et un journal de quinze décisions
représentent un investissement documentaire significatif, supérieur à ce
qu’un développement humain classique consent — l’étude de cas l’a payé
parce qu’il s’agissait précisément de valider la discipline. Elle a un
coût de tempo~: aucune écriture n’est permanente sans vérification
complète, ce qui rend les itérations plus lentes qu’en édition libre — le
prix payé pour que la régression soit impossible \emph{par construction}
plutôt que rare en pratique. Elle a enfin un coût d’infrastructure~:
intégration continue multi-OS, triple sauvegarde, reproducteur
d’empaquetage. L’argument économique du protocole est que ces coûts sont
des \emph{investissements anti-dette}~: dans l’étude de cas, le diagnostic
complet du seuil D8 — mesure, décision, spécification, tests — a coûté une
fraction de ce qu’aurait coûté le seul dépannage du défaut en production~;
et le défaut existait depuis trois ans parce que rien ne l’avait rendu
visible. Ces observations restent qualitatives~; aucune mesure
comparative du coût total n’a été tentée (section~\ref{sec:menaces}).

\subsection{Enseignements pour le protocole lui-même}

L’étude renvoie au protocole trois enseignements concrets, que la
prochaine révision pourra instruire. Premièrement, la notion de
\emph{pré-version de cristallisation} (\code{-dev.N})~: la démonstration a
distingué « vérifié par l’oracle » de « validé sur le terrain », nuance
que la spécification ne prévoyait pas. Deuxièmement, la \emph{reverse-spec}
elle-même, qui étend la couche I aux binaires scellés hérités. Troisièmement,
la valeur du \emph{partage de fixtures} entre consommateurs d’un même
contrat externe — les spécifications de l’installateur et l’application
partagent contractuellement leurs jeux de test d’analyse de sortie, ce qui
élimine par construction la duplication qui laisse vivre les défauts en
trois exemplaires. Chacun de ces enseignements est entré au journal des
décisions du dépôt~; aucun n’a été appliqué en silence — la gouvernance
s’applique à elle-même.
